\documentclass[10pt,a4paper]{amsart}
\usepackage[margin=19mm]{geometry}
\usepackage{amsmath,amssymb,mathtools}
\usepackage[scr=rsfs,cal=euler]{mathalfa}
\usepackage{xfrac}
\usepackage[unicode,hidelinks,bookmarks=false]{hyperref}
\newcommand{\ie}{i.e.\ }
\newcommand{\cf}{cf.\ }
\newcommand{\resp}{respectively\ }
\newcommand{\Subsection}{\subsection}
\providecommand{\nobreakdash}{\nobreak\hbox{-}}
\setlength{\emergencystretch}{2em}
\multlinegap0pt
\usepackage{amssymb}
\usepackage{mathtools}
\usepackage{bm}
\usepackage[shortlabels]{enumitem}
\newtheorem{assumption}{Assumption}
\newtheorem{theorem}{Theorem}
\newcommand{\C}{\mathbb C}
\newcommand{\Ga}{\Gamma}
\newcommand{\norm}[1]{\left\|#1\right\|}
\title{Periodic B-spline--Heaviside Collocation for Fredholm Integral Equations with Piecewise H\"older Data}
\author{Maria Capcelea, Titu Capcelea}
\date{Source: arXiv:2511.06590v2 (CC BY 4.0)}
\begin{document}
\maketitle

\noindent\emph{Source and licence.} Periodic B-spline--Heaviside Collocation for Fredholm Integral Equations with Piecewise H\"older Data. Version arXiv:2511.06590v2 (CC BY 4.0), distributed by arXiv under the Creative Commons Attribution 4.0 International licence, http://creativecommons.org/licenses/by/4.0/, as recorded in the arXiv licence metadata for this exact version and shown on the abstract page https://arxiv.org/abs/2511.06590v2. Full licence notice: \texttt{licenses/arxiv-2511.06590-CC-BY-4.0.txt}.\par\smallskip

\noindent\textit{Editorial scope.} Counted unit: the article's theorem thm:bsh-iterated (source0.tex lines 893--914). The article's complete original proof of it is delivered with it (source0.tex lines 916--955, one \texttt{proof} environment of 40 lines). No mathematics is written, repaired or re-derived: the statement and the proof are the article's own lines, in the article's own notation, and no retained line is substituted or re-flowed.  Retained uncounted as the source-local context the proof needs: lines 217--233; lines 266--276; lines 522--525; lines 613--623; lines 888--891.  Nothing was omitted from any retained span, so the package carries no omission record.  No retained line contains a citation, so the wrapper carries no bibliography and no bibliography entry.  No retained line contains a figure environment, an image-inclusion command or drawing code, so no image is delivered and the package declares no asset dependency.  Editorial formatting only, in addition: an AMS \texttt{amsart} preamble that restates the article's own declarations and macros used by the retained lines, namely the environments \texttt{assumption}, \texttt{theorem} on separate counters, together with the standard packages those lines need.  The retained environments print in this document as Assumption 1, Theorem 1 in order of appearance, whereas the article prints the same items with the numbering of its own full document.  The complete native source of the article as distributed in the arXiv e-print for this exact version is bundled unchanged as \texttt{sources/188744/source0.tex}.
\begin{assumption}[Kernel regularity]\label{ass:kernel}
There exists an exponent $\mu$ such that
\[
        0<\beta<\alpha<\mu\le1
\]
and constants $M_0,M_\mu$ for which
\begin{align}
        |k(t,\tau)|&\le M_0,\label{eq:kbound}\\
        |k(t,\tau)-k(s,\tau)|&\le M_\mu|t-s|^\mu\label{eq:kholder}
\end{align}
for all $t,s,\tau\in\Ga$.  In addition, the family $\{\tau\mapsto k(t,\tau):t\in\Ga\}$ is equicontinuous.  Equivalently, there is a modulus of continuity $\omega_k$, with $\omega_k(r)\to0$ as $r\downarrow0$, such that
\begin{equation}\label{eq:kuniformtau}
        |k(t,\tau)-k(t,\zeta)|
        \le \omega_k(|\tau-\zeta|),
        \qquad t,\tau,\zeta\in\Ga.
\end{equation}
\end{assumption}

\begin{assumption}[Nonresonance of the parameter]\label{ass:injectivity}
The fixed parameter $\lambda\in\C$ is such that
\begin{equation}\label{eq:nonresonance}
        \ker(I-\lambda K)=\{0\}
        \qquad\text{in }X_\beta.
\end{equation}
Equivalently, either $\lambda=0$, or
\[
        \lambda^{-1}\notin\sigma(K:X_\beta\to X_\beta).
\]
\end{assumption}

\begin{equation}\label{eq:ubound}
        \norm{u}_{H_\alpha(\Ga)}
        \le C\norm{f}_{X_\alpha}.
\end{equation}

\begin{align}
        \norm{P_hv}_{\infty}
        &\le C\norm{v}_{\infty},\label{eq:Phsupstable}\\
        \norm{v-P_hv}_{\infty}
        &\le Ch^\alpha\norm{v}_{H_\alpha(\Ga)},\label{eq:Phsupapprox}\\
        \norm{P_hv}_{H_\beta(\Ga)}
        &\le C\norm{v}_{H_\beta(\Ga)},\label{eq:Phstable}\\
        \norm{v-P_hv}_{H_\beta(\Ga)}
        &\le Ch^{\alpha-\beta}\norm{v}_{H_\alpha(\Ga)}.
        \label{eq:Phapprox}
\end{align}

\begin{equation}\label{eq:bsh-iterated-error-identity}
        \varphi-\varphi_h^{\mathrm{it}}
        =\lambda K(\varphi-\varphi_h).
\end{equation}

\begin{theorem}[Improved raw and iterated estimates]
\label{thm:bsh-iterated}
Let
\[
        0<\beta<\alpha<\mu\le1,
\]
and suppose Assumptions~\ref{ass:kernel} and~\ref{ass:injectivity} hold.
Then, for all sufficiently small $h$, the exact B-spline--Heaviside
collocation solution satisfies
\begin{equation}\label{eq:bsh-raw-sup}
        \norm{\varphi-\varphi_h}_{\infty}
        \le
        C_\lambda h^\alpha\norm{f}_{X_\alpha},
\end{equation}
and its iterated extension satisfies
\begin{equation}\label{eq:bsh-iterated-bound}
        \norm{\varphi-\varphi_h^{\mathrm{it}}}_{H_\mu(\Ga)}
        \le
        C_\lambda h^\alpha\norm{f}_{X_\alpha}.
\end{equation}
Consequently, \eqref{eq:bsh-iterated-bound} also holds in $C(\Ga)$ and in every $H_\sigma(\Ga)$ with $0<\sigma\le\mu$.
\end{theorem}

\begin{proof}
First consider $A=I-\lambda K$ on $C(\Ga)$ with the uniform norm.
The map $K:C(\Ga)\to H_\mu(\Ga)$ is bounded and compact as a map into
$C(\Ga)$.  If $Av=0$ for some $v\in C(\Ga)$, then
$v=\lambda Kv\in H_\mu(\Ga)\subset X_\beta$, and Assumption
\ref{ass:injectivity} gives $v=0$.  Hence $A$ is boundedly invertible on
$C(\Ga)$.

By \eqref{eq:Phsupstable}, \eqref{eq:Phsupapprox}, and the smoothing property
of $K$,
\[
        \norm{(I-P_h)K}_{C(\Ga)\to C(\Ga)}
        \le Ch^\mu.
\]
Therefore $A_h=I-\lambda P_hK$ is uniformly invertible on $C(\Ga)$ for all
sufficiently small $h$.  Since
\[
        A_h(u-u_h)=(I-P_h)u,
\]
we obtain from \eqref{eq:Phsupapprox} and \eqref{eq:ubound}
\[
        \norm{u-u_h}_{\infty}
        \le
        C_\lambda h^\alpha\norm{u}_{H_\alpha}
        \le
        C_\lambda h^\alpha\norm{f}_{X_\alpha}.
\]
The exact and discrete jump components coincide, so
$\varphi-\varphi_h=u-u_h$, which proves \eqref{eq:bsh-raw-sup}.

Finally, \eqref{eq:bsh-iterated-error-identity} and the bound
$K:C(\Ga)\to H_\mu(\Ga)$ give
\[
        \norm{\varphi-\varphi_h^{\mathrm{it}}}_{H_\mu}
        \le
        |\lambda|\,C\norm{\varphi-\varphi_h}_{\infty}.
\]
Combining this estimate with \eqref{eq:bsh-raw-sup} proves
\eqref{eq:bsh-iterated-bound}.
\end{proof}

\end{document}
